\section{Related work}
\label{sec:related}

\subsection{Snow-depth and SWE mapping from terrain and remote sensing}

The spatial distribution of mountain snow has traditionally been modelled by
relating snow depth or snow water equivalent (SWE) to terrain attributes
through statistical regression. Elevation, slope, aspect (often decomposed
into northness and eastness), curvature and topographic position indices are
the classical predictors \citep{elder1991,molotch2005,revuelto2014}. Wind
redistribution, a dominant control in alpine terrain, is commonly represented
by terrain-based parameters such as the maximum upwind slope $S_x$
\citep{winstral2002,winstral2013}, which we also adopt here. More recently,
machine-learning regressors -- particularly Random Forests
\citep{breiman2001} -- have become a de-facto standard for HS/SWE
interpolation because they handle non-linear interactions and large predictor
sets robustly \citep{broxton2019,king2020,ntokas2021}. These pixel-wise
methods are strong baselines but, by construction, ignore the spatial
arrangement of neighbouring pixels.

\subsection{Deep learning for high-resolution environmental mapping}

Convolutional neural networks have transformed dense prediction tasks in
remote sensing, from land-cover classification and change detection
\citep{daudt2018} to super-resolution and regression of continuous
geophysical fields. Encoder--decoder architectures with skip connections, of
which the U-Net \citep{ronneberger2015} is the archetype, are especially
well suited to mapping problems because they combine multi-scale context with
spatially precise outputs. Residual and attention-augmented variants such as
ResUNet \citep{zhang2018} and ResUNet++ \citep{jha2019} add residual blocks,
squeeze-and-excitation and atrous spatial pyramid pooling, which stabilise
training and improve the modelling of multi-scale structures. In the snow
context, CNNs and other deep models have been explored for snow-cover
mapping, SWE reconstruction and downscaling
\citep{king2020,mital2022}, but high-resolution (metre-scale) HS regression
from LiDAR remains comparatively under-explored, and systematic comparisons
against strong tabular baselines under temporally disjoint evaluation are
scarce. The work most closely related to ours is MAPunet \citep{betato2025mapunet},
which casts high-resolution HS mapping as U-Net \emph{pixel-wise regression}
from a multi-channel stack of topographic and snow-related predictors and
applies it over the Swiss Alps (Davos). The present study builds directly on
that pixel-wise regression framework but (i) replaces the plain U-Net backbone
with a residual, attention-augmented ResUNet++; (ii) adds a pattern-oriented
spatial loss and pattern-aware metrics (\SPAEF{}/\MSPAEF{}); and (iii) evaluates
under a strictly temporal split with a multi-seed stability analysis.

\subsection{Evaluation metrics for spatial fields}

Pixel-wise metrics (\Rsq{}, RMSE, MAE, bias, Nash--Sutcliffe efficiency)
dominate the snow-modelling literature, yet they assess only the marginal
agreement between predicted and observed values and are blind to the spatial
organisation of the error. The Spatial Efficiency metric (\SPAEF{})
introduced by \citet{koch2018} was designed precisely to evaluate the fidelity
of spatial patterns, by combining three complementary components: the Pearson
correlation between the two fields, the ratio of their coefficients of
variation, and the overlap of their value histograms. \SPAEF{} and related
multi-component metrics have since been adopted to benchmark distributed
hydrological and land-surface models \citep{koch2018,demirel2018}. Despite
their suitability, such pattern-oriented metrics are rarely used for
high-resolution snow mapping. We argue, and demonstrate empirically, that they
are necessary to obtain an honest picture of model quality and stability, and
we further use them -- through a differentiable spatial-correlation term -- to
\emph{guide} the optimisation, not merely to evaluate it.

\subsection{Pattern-oriented and physics-aware losses}

Optimising solely a pixel-wise loss (MSE/MAE) tends to produce regression to
the mean and over-smoothed fields. To counteract this, a growing body of work
augments the training objective with structural or perceptual terms -- for
example structural similarity (SSIM) losses, gradient-difference losses, or
adversarial terms -- so that the model is rewarded for reproducing spatial
structure \citep{zhao2017}. In the environmental domain, correlation-based and
multi-scale loss terms have been used to better preserve spatial texture.
Building on this idea, we propose a simple and interpretable hybrid loss that
linearly blends MSE with a spatial Pearson-correlation penalty controlled by a
single hyperparameter $\lpear$, and we characterise its effect on both
accuracy and spatial pattern fidelity through a full sweep.
